\lesson{1}{Mar 18 2022 Fri (16:40:24)}{Write an inequality from Graph}{Unit 2}
\label{les_1_unit_2:write_an_inequality_from_graph}

To find the inequality of a graph is very similar to finding the equation to a
linear graph. You would just find out the $y$-intercept and the slope
(\ref{def:slope}). The difference is that if the line is dashed, then instead of
an $=$ symbol, it's going to be either $<$ or $>$. But, if it's a straight line,
you would use $\le$ or $\ge$ in replace for the $=$ symbol.

The $>$ and $\ge$ means the shaded area is below the line because it's saying
that the $y$ value is bigger than the $x$ value.

The $<$ and $\le$ means the shaded area is above the line because it's saying
that the $y$ value is less than the $x$ value.

\begin{example}
  \label{exm:write_an_inequality_from_graph}

  Let's say we have the following inequality graph:

  \begin{figure}[H]
    \centering
    \incfig{write-an-inequality-from-graph}
    \label{fig:write_an_inequality_from_graph}
  \end{figure}

  We already know that $b$ is $2$, because that's the $y$-intercept. To find the
  slope, let's just pick two random points. I'm going with:
  $(0,2) \qquad (1, 6)$. Let's use the slope (\ref{def:slope}) formula to find
  the slope of this line.

  \begin{equation}
    \label{eqt:using_slope_formula_to_find_slope_of_inequality_line}
    \begin{split}
      \frac{y_2 - y_1}{x_2 - x_1} &= \frac{6 - 2}{1 - 0} \\
                                  &= \frac{4}{1} \\
                                  &= 4
    \end{split}
  .\end{equation}

  So, $m = 4$, which is our slope. So, here's our final equation: 

  \[ y = 4x + 2 \].

  To find the inequality symbol, we look at the shaded area. Since the shaded
  area is below the line, then that means the symbol will either be: $<$ or
  $\le$. Since the line is straight (meaning, not dashed), that means that it
  also includes all of the numbers that the line is on. For that, we use the
  $\le$ symbol.

  Here's our final answer:

  \[ \boxed{y \le 4x + 2} \].
\end{example}

\newpage
